\noindent\textbf{(a)} $y=(C/x+x/2)e^{-2x}$ para $x\ne0$. Sólo $C=0$ permite la prolongación por $x=0$, que da $y=(x/2)e^{-2x}$ en toda la recta.

\noindent\textbf{(b)} $y=(x+1)^2(1+x+x^2/2)$, con intervalo máximo $(-1,\infty)$.

\noindent\textbf{(c)} $y=2-(1+x^2)^{-3/2}$ en $\mathbb R$.

\noindent\textbf{(d)} $y=(x^3-3x+C)/[3x(x+1)]$ para $x\ne0,-1$. Si $C=0$, se prolonga por $x=0$ como $y=(x^2-3)/[3(x+1)]$, con $y(0)=-1$. Siempre se excluye $x=-1$.
